\begin{problem}[Warmup]{(a)}
    Compute a quotient and remainder for the following pairs of elements, in the given Euclidean domain.
    \begin{enumerate}
        \item[(a)] 116 divided by 7 in $\mathbb{Z}$.
        \item[(b)] $x^4 - 2x^3 + 5x - 3$ divided by $2x^2 - 1$ in $\mathbb{Q}[x]$.
        \item[(c)] $13 + 5i$ divided by $3 + 2i$ in $\mathbb{Z}[i]$.
        \item[(d)] 7 divided by $\pi$ in $\mathbb{R}$.
    \end{enumerate}
\end{problem}
\unsolved

\begin{problem}[Warmup]{(b)}
    Given an ideal $I$ in a Euclidean domain, we know that $I$ is principal. Describe conceptually how to determine the generator of $I$. (This description is very simple; it might not be so simple to carry out in practice.)
\end{problem}
\unsolved

\begin{problem}{1}
    Using a similar argument to that in the text for $\mathbb{Z}[i]$, prove that
    $\mathbb{Z}[\sqrt{2}]$ is a Euclidean domain, using the function $N$ as the valuation.
\end{problem}
\unsolved

\begin{problem}{2}
    Compute a quotient and remainder for $17 + 32\sqrt{2}$ divided by $3 - 4\sqrt{2}$ in
    the Euclidean domain $\mathbb{Z}[\sqrt{2}]$. Check that your remainder has a smaller
    valuation than $3 - 4\sqrt{2}$.
\end{problem}
\unsolved

\begin{problem}{3}
    Show by example that quotients and remainders are not unique in the Euclidean domain
    $\mathbb{Z}[\sqrt{2}]$.
\end{problem}
\unsolved

\begin{problem}{4}
    Make a definition for \emph{greatest common divisor} of two elements in a Euclidean
    domain.
\end{problem}
\unsolved

\begin{problem}{5}
    Suppose that $a, b \in D$, and $D$ is a Euclidean domain. Consider the ideal
    $\langle a, b \rangle = \{ax + by : x, y \in D\}$; see Exercise 12.1. Let $d$ be an
    element of $\langle a, b \rangle$ with least valuation.
    \begin{enumerate}
        \item[(a)] Why does the element $d$ exist?
        \item[(b)] Prove that $d$ is a gcd of $a, b$. (See Exercise 4.)
        \item[(c)] Why does this mean that Euclidean domains possess a GCD identity?
        \item[(d)] Rephrase the GCD identity for Euclidean domains in terms of ideals of
        the form $\langle a, b \rangle$.
    \end{enumerate}
\end{problem}
\unsolved

\begin{problem}{6}
    \begin{enumerate}
        \item[(a)] Explain why Euclid's Algorithm for finding a gcd makes sense in a
        Euclidean domain.
        \item[(b)] Apply Euclid's Algorithm to $5 + 133i$ and $17 + 34i$ in $\mathbb{Z}[i]$.
        \item[(c)] Backtrack through Euclid's Algorithm to show explicitly that
        $\langle 5 + 133i, 17 + 34i \rangle$ is a principal ideal.
    \end{enumerate}
\end{problem}
\unsolved

\begin{problem}{7}
    Suppose that $a$ and $b$ are integers. Prove that their gcd in $\mathbb{Z}[i]$ is the
    same as in $\mathbb{Z}$, up to unit multiples in $\mathbb{Z}[i]$.
\end{problem}
\unsolved

\begin{problem}{8}
    Prove that if $D$ is a Euclidean domain and $a, b \in D$ are associates, then
    $v(a) = v(b)$. (This is easy!) Show by example that the converse is false in
    $\mathbb{Z}[\sqrt{2}]$. (You should compare this exercise to Exercise 10.14, where two
    non-associates of the same norm are exhibited in $\mathbb{Z}[\sqrt{-5}]$; of course,
    the latter ring is not a Euclidean domain.)
\end{problem}
\unsolved

\begin{problem}{9}
    Provide a geometric interpretation of the proof that $\mathbb{Z}[\sqrt{2}]$ is a
    Euclidean domain, analogous to the interpretation we discussed for $\mathbb{Z}[i]$.
\end{problem}
\unsolved
